\section{Quantitative details for the corner maximum}
\label{app:argmax-details}

This appendix records the quantitative comparisons underlying
Lemma~\ref{lem:argmax-four-regions}.  The constants are sufficient
constants; no minimal analytic onset is asserted.  All logarithms are
natural.  Statements about real $n$ below concern explicit scalar
allowances, whereas interpolation objects always have integer degree.

\subsection{Exact reflection identities}

Write $a_j=j\pi/n$, $b_k=k\pi/m$, and let $e_j,f_k$ be one in the
interior of their respective index intervals and $1/2$ at either
endpoint.  The cardinal weights are $w_{jk}=2e_jf_k/\mathsf N$.  Product to sum
in the corrected cardinal formula gives
\[
 L_{jk}(\cos\theta,\cos\phi)
 =\frac{w_{jk}}4\sum_{\sigma,\eta=\pm1}
       G_n(\theta+\sigma a_j,\phi+\eta b_k).
\]
For $g_{\sigma\eta}=G_n(\theta+\sigma a_j,\phi+\eta b_k)$ set
\begin{align*}
 M(\theta,\phi)&=\sum_{j+k\ \text{even}}
                     \frac{w_{jk}}4\sum_{\sigma,\eta}|g_{\sigma\eta}|,\\
 \Delta(\theta,\phi)&=\sum_{j+k\ \text{even}}\frac{w_{jk}}4
       \left(\sum_{\sigma,\eta}|g_{\sigma\eta}|
                       -\left|\sum_{\sigma,\eta}g_{\sigma\eta}\right|\right).
\end{align*}
Then $\lambda_n=M-\Delta$ exactly.  For every residue $r$ modulo $2n$,
\[
 \sum_{\substack{0\leq j\leq n,\ \sigma=\pm1\\
                   \sigma j\equiv r\ (2n)}}e_j=1.
\]
The same identity holds for the $m$ mesh.  Consequently
\[
 M(\theta,\phi)=\frac1{2\mathsf N}
 \sum_{\substack{0\leq r<2n,\ 0\leq s<2m\\r+s\ \text{even}}}
       |G_n(\theta+r\pi/n,\phi+s\pi/m)|.
\]
Translation by an odd center $(J,K)$ turns this even board into the odd
board in \eqref{eq:argmax-primitive}.  This proves the exact identity
\eqref{eq:argmax-defect}, without an endpoint correction.

For completeness, put $d(v,w)=|v|+|w|-|v+w|$,
$\Gamma^-=\{(r,s):0\leq r<2n,\ 0\leq s<2m,\ r+s\text{ odd}\}$, and
\[
 H_{rs}=G_n\left(\frac{r\pi+\alpha}{n},
                    \frac{s\pi+\beta}{m}\right),\qquad (r,s)\in\Gamma^-.
\]
With periodic indices,
\begin{equation}
 Q_v(K)=\frac1{4\mathsf N}\sum_{\Gamma^-}d(H_{rs},H_{r,2K-s}),\qquad
 Q_h(J)=\frac1{4\mathsf N}\sum_{\Gamma^-}d(H_{rs},H_{2J-r,s}).
                                                        \label{eq:argmax-projections}
\end{equation}
The phases are unchanged in the reflected partners.  Applying the
triangle inequality first within either pair gives
$\Delta\geq Q_v$ and $\Delta\geq Q_h$.  These are separate projections;
their sum is not a lower bound used here.

All formulas are finite-polynomial identities.  For evaluation one may
use $P_n=T_{n+1}+T_n$ and
\[
 G_n(u,v)=\frac{P_n(\cos u)-P_n(\cos v)}{\cos u-\cos v}-\cos(nu)
\]
away from a zero denominator, and replace the quotient by
$P_n'(\cos u)$ at equality.  No collision term is discarded.  In
particular $G_n(0,0)=2\mathsf N$, $G_n(\pi,\pi)=2n(-1)^n$, and
$G_n(0,\pi)=0$.

\subsection{The paired central contribution}

The completed model and its cancellation functional are developed in
the supporting model estimates.  We give the joint argument responsible
for the useful sign of the second phase.  Put
\begin{gather*}
 x=|\alpha|,\quad y=|\beta|,\quad
 a=\frac{\cos x+\cos y}{2},\quad u=1-a,\\
 p=\sin x,\quad q=\sin y,\quad v=p/2,\quad\delta=q/2.
\end{gather*}
Let $J(v,s,e)$ be the generic cancellation integral in
Lemma~\ref{lem:generic-cancellation}, and define
\begin{align}
 s_0(e)&=a\tan(e/2),\qquad
 s_L(e)=s_0(e)+\delta,\qquad s_R(e)=|s_0(e)-\delta|,\nonumber\\
 E_c(x,y)&=\frac1{\pi^2}\int_0^{\pi/2}
 \{J(v,s_L(e),e)+J(v,s_R(e),e)-4s_0(e)\log\sin e\}\,de.
                                                        \label{eq:argmax-central-integral}
\end{align}
Write $E(v,s,e)=J(v,s,e)-2s\log\sin e$ and
\[
 Z(v,s,e)=2\int_e^{\pi/2}
           \frac{\sqrt{(v^2-s^2\sin^2t)_+}}{\sin t}\,dt.
\]
The generic cancellation bound is $E\leq Z$.  Since
$s_L+s_R-2s_0=2(\delta-s_0)_+$, an exact rearrangement gives
\begin{align}
 E_c&=\frac1{\pi^2}\int_0^{\pi/2}
              \{E(v,s_L(e),e)+E(v,s_R(e),e)\}\,de-B_\beta,\nonumber\\
 B_\beta&=\frac4{\pi^2}\int_0^b
        \{q/2-a\tan(e/2)\}[-\log\sin e]\,de,\qquad
 b=2\arctan\frac q{2a}.                             \label{eq:argmax-beta-baseline}
\end{align}
Here $b=B_\beta=0$ if $q=0$.  On $x+y\leq1$, $a>0$ and
$0\leq b\leq x+y<\pi/2$, so $B_\beta\geq0$.

\begin{lemma}[Paired root bound]
\label{lem:argmax-central-pair}
Let $I_4=\int_0^1\sqrt{1-t^4}\,dt$ and $\kappa_4=4I_4/\pi^2$.
On $x,y\geq0$, $x+y\leq1$,
\begin{equation}
 E_c(x,y)\leq\kappa_4\frac{p^{3/2}}{\sqrt a}-B_\beta,
 \qquad
 E_c\leq
 \begin{cases}
 3x^{3/2}/8-B_\beta,&x+y\leq1/2,\\
 5x^{3/2}/12-B_\beta,&x+y\leq1.
 \end{cases}                                         \label{eq:argmax-central-pair}
\end{equation}
The exact boundary identity is $E_c(0,y)=-B_\beta$.
\end{lemma}

\begin{proof}
After applying $E\leq Z$, it suffices to bound
\[
 \mathcal R=\int_0^{\pi/2}
       \{Z(v,a\tan(e/2)+\delta,e)+Z(v,|a\tan(e/2)-\delta|,e)\}\,de.
\]
For fixed $0<t\leq\pi/2$ define on the real line
\[
 W_t(w)=\frac{2a}{a^2+w^2}{\bf1}_{|w|\leq a\tan(t/2)},
 \qquad R_t(w)=\sqrt{(v^2-w^2\sin^2t)_+}.
\]
Tonelli's theorem and $w=a\tan(e/2)$ give exactly
\[
 \mathcal R=2\int_0^{\pi/2}\frac1{\sin t}
              \int_{\mathbb R}W_t(w)R_t(w-\delta)\,dw\,dt.
\]
Both functions in the inner convolution are nonnegative, even, and
nonincreasing with $|w|$.  The convolution is largest at zero: by
layer cake, its value is the integral of intersection lengths of two
intervals, one translated by $\delta$; each intersection has length at
most that of the smaller interval, attained when both are centered.
Thus
\[
 \mathcal R\leq2\int_0^{\pi/2}Z(v,a\tan(e/2),e)\,de.
\]
There is no signed interchange here, and
$Z(v,s,e)\leq2v\log(2/e)$ ensures finiteness.  The root integral bound
of Lemma~\ref{lem:root-majorant} is
\[
 \int_0^{\pi/2}Z(v,a\tan(e/2),e)\,de
       \leq4\sqrt2 I_4\frac{v^{3/2}}{\sqrt a},\qquad 0\leq v\leq a.
\]
On the stated domain, $a\geq49/64$ and $v/a\leq32/49<1$.
Substitution of $v=p/2$ proves the first inequality, including its
normalization.  The elementary conversions in the root-majorant lemma,
using $I_4<319/360$, give the two rational constants.  At $v=0$ the
generic cancellation identity is exact with $E=0$, proving the boundary
identity.
\end{proof}

On the smaller diamond the negative baseline is quantitatively useful.
Indeed $2a\leq1+\cos y$ gives $b\geq y$, while $b\leq x+y\leq1/2$.
Convexity of tangent yields
$a\tan(e/2)\leq qe/(2b)$ for $0\leq e\leq b$.  Since $\sin e\leq e$,
\begin{equation}
 B_\beta\geq\frac{qb}{\pi^2}\left(\log\frac1b+\frac32\right)
       \geq\frac{1679}{8000}y^2.                   \label{eq:argmax-beta-quadratic}
\end{equation}
For the last comparison use $q\geq23y/24$, $b\geq y$, $\pi^2<10$,
and $\log2>69/100$.  This also holds at $y=0$ by continuity.

\subsection{Scalar allowances and their roles}

Set $L=\log n$, $H=1+L$, $b_0=4/\pi^2$, and
$h_r=\log(5/2)+1/18$.  The explicit allowances used below are
\begin{align}
 T&=\frac{4L^2-3L}{\pi^2},\nonumber\\
 D&=\frac{(52/\pi^2)H+7}{n}+\frac{10}{n^2}
                                      +\frac{20H}{\pi n^2},\nonumber\\
 \mathcal M&=\frac{b_0(2L+2h_r)}n
 +\frac{4+b_0(L+1/n)}{4n^2-1}+\frac{b_0}{(2n-1)^2}\nonumber\\
 &\hspace{12mm}+\frac{5+b_0(\log(n+1/2)+2)}{4(n-1/2)^2},\nonumber\\
 \bar e&=D+\mathcal M+1/n,\qquad
 K_D=\frac{[6+16/(3\pi)]H+2/\pi}{n},\nonumber\\
 A_4&=4/n+40H/(\pi n^2),\qquad C_4=A_4+20/n^2,\nonumber\\
 d_\beta&=(17L+27)/n^2,\nonumber\\
 J_N&=\frac{(3/2)H+16/5}{n}+\frac{78}{n^2},\nonumber\\
 J_W&=\frac{(3/2)H+6}{n}+\frac{148}{n^2},\nonumber\\
 k_{Nx}&=\frac{(19/2)H+(23/3)L+677/10}{n}+\frac{82}{n^2},
 \quad k_{Ny}=k_{Nx}+d_\beta,\nonumber\\
 k_{Wx}&=\frac{(19/2)H+(23/3)L+155/2}{n}+\frac{152}{n^2},
 \quad k_{Wy}=k_{Wx}+d_\beta,\nonumber\\
 U_N&=T-\bar e+A_4+\rho J_N+j,\nonumber\\
 U_W&=T-\bar e+A_4+\rho J_W+j.                    \label{eq:argmax-coefficients}
\end{align}
The letter $\mathcal M$ here denotes a scalar mesh remainder, not a
global norm.  The regular-model estimates in
Proposition~\ref{scalar:origin-bounds} and
\eqref{scalar:global-nine-fourths} supply
\begin{equation}
 j>23/250,\qquad t_0=2\log2/\pi<221/500,
 \qquad C_\infty(\alpha,\beta)<9/4.                 \label{eq:argmax-regular-constants}
\end{equation}
The origin lemma gives only $E_n(0)>-\bar e$ from $n=160$;
it is not a global absolute bound.  The exact relative conversion is
\[
 E_n(\alpha,\beta)-E_n(0)
 =\{E_n(\alpha,\beta)-aE_n(0)\}-uE_n(0).
\]
The complete row estimates therefore give
\begin{equation}
 S_n-\Lambda_c\leq
 \begin{cases}
 E_c-uU_N+C_4+k_{Nx}x+k_{Ny}y,&\rho\leq1/2,\\
 E_c-uU_W+C_4+k_{Wx}x+k_{Wy}y,&\rho\leq1.
 \end{cases}                                         \label{eq:argmax-relative-comparison}
\end{equation}
The second line uses the wide signed outer bound, not an extension of
the sharper half-diamond bound.  The separate mesh charge is
$\rho/(2n)$; the scalar logarithmic term is already inside the signed
outer estimate.  Lemma~\ref{lem:relative-telescope} gives the complete
component matching.

For the corner argument we also require the full absolute allowance.
Define
\begin{align*}
 g&=1+\log(2+8m/(3\pi)),\qquad M_A=1/\pi+4/\pi^2,\\
 G_g&=\frac1m\left\{7M_A L+\frac{111}2+
 \frac{63+\log2+4\log\pi+7\log m+16\log(2m)+(3/2)g}{\pi}\right\},\\
 W&=\frac{3\log m+\log2+2\log\pi+2}{\pi m},\\
 \overline C_E&=D+G_g+(64+16L)/n+318/n^2+\mathcal M+W.
\end{align*}
The full-phase theorem, Lemma~\ref{lem:phase-row-global}, gives
\begin{equation}
 |S_n-M_\infty|\leq\overline C_E+K_D P,
 \qquad P=\sin x+\sin y\leq2,\qquad \rho\leq\pi.
                                                        \label{eq:argmax-global-error-norm}
\end{equation}
Combining it with the global model coefficients and
\eqref{eq:argmax-regular-constants} yields the genuine global norm
\begin{equation}
 M(n)=\frac{4L^2+(79/5)L}{\pi^2}+\frac94+\overline C_E+2K_D.
                                                        \label{eq:argmax-analytic-norm}
\end{equation}
Even-lattice phase translations and separate phase reflections reduce
all real phases to $\rho\leq\pi$.  Thus \eqref{eq:argmax-analytic-norm}
bounds every complete signed board sum.  By the exact reflection
identity it also bounds every complete signed cardinal sum after the
inward corner substitution, globally in its real phase variables.

The quantitative trust in \eqref{eq:argmax-regular-constants} must be
retained.  With $k_s=\gamma+\log(8/\pi)$, its origin normalization uses
\[
 C_* =4k_s^2/\pi^2+\kappa_c,
\]
\begin{gather*}
 0.919710118388492535099518225790<\kappa_c,\\
 \kappa_c<0.919710118388492535099518225791,\\
 1.846160310423926079759838737947<C_*,\\
 C_*<1.846160310423926079759838737948.
\end{gather*}
These are the inherited PD-T271 outward enclosures, obtained from a
coefficient series through index $64$ at $320$-bit Arb precision and
the analytic tail radius $8/(3\cdot66\cdot4^{65})$.  They concern a
limiting constant.  No finite lower bound for $\Lambda_c(n)$ is inferred
from these decimal endpoints alone.  The analytic scalar maximum
argument does not remove this numerical dependency.

\subsection{The half-near comparison from degree 320}
\label{app:argmax-half-near}

The elementary identity
$\log z=2\sum_{r\geq0}((z-1)/(z+1))^{2r+1}/(2r+1)$, with a
geometric bound on its positive tail, gives at $n=320=2^8(5/4)$
\[
 \frac{57679}{10000}<L<\frac{5769}{1000},\qquad H<\frac{6769}{1000}.
\]
For explicit rational brackets, $\log2$ lies between
$842/1215+2/15309$ and that sum plus $1/78732$; $\log(5/4)$ lies
between $488/2187$ and that sum plus $1/145800$.
Together with $157/50<\pi<22/7$, $\pi^2>197/20$, and $h_r<1$,
substitution into \eqref{eq:argmax-coefficients} yields
\begin{equation}
 \begin{gathered}
 T>293/25,\quad D<27/200,\quad\mathcal M<7/400,
 \quad\bar e<249/1600,\quad C_4<677/50000,\\
 k_{Nx}<11031/20000,\qquad k_{Ny}<11056/20000,
 \qquad J_N>2171/51200>21/500,\\
 U_N>23337/2000+(21/500)\rho.
 \end{gathered}                                      \label{eq:argmax-base320}
\end{equation}
For example the lower numerator for $T$ is
$567277810036/48400000000>293/25$.  For $D$ a valid upper bound is
\[
 \frac{(1040/197)(6769/1000)+7}{320}
       +\frac{10+6769/157}{320^2}<\frac{27}{200}.
\]
The leading mesh term is at most
$(41/100)(2(5769/1000)+2)/320$; the other three positive terms total
less than $14/320^2$, proving the displayed mesh bound.  The exact
upper numerator for $k_{Nx}$ is $705963/1280000$; addition of the
entire $d_\beta$ gives $56602113/102400000$ for $k_{Ny}$.

These bounds persist with the needed directions for all real $n\geq320$.
$\bar e,C_4,k_{Nx},k_{Ny}$ decrease.  For the whole mixed coefficient
$U_N$, use $H\leq n/4$ on $n\geq32$, $L>5$, and $\pi^2<10$:
\[
 T'>3/n,\quad A_4'>-32/(3n^2),\quad
 J_N'>-3/(8n)-16/(5n^2)-156/n^3.
\]
For $0\leq\rho\leq1/2$ this implies
\[
 (T+A_4+\rho J_N)'
 >\frac{45}{16n}-\frac{184}{15n^2}-\frac{78}{n^3}>0.
\]
The decreasing $\bar e$ improves the direction.  This checks the
complete coefficient, rather than assuming a favorable direction for
a difference of positive terms.

The projected-floor formulas in Lemma~\ref{lem:projected-floors} give,
for each eligible interior-index projection,
\begin{equation}
 \Delta_{JK}>L_h(\rho):=
 \frac{5123}{50000}-\frac{4447}{50000}\rho
                         -\frac{1281}{50000}\rho^2,
 \qquad 0\leq\rho\leq\frac12.                      \label{eq:argmax-half-floor}
\end{equation}
Here is the base calculation, including every center class.  Write
$a_\rho=1-\rho^2/4$.  In the middle class the formula is
$F_h=A_h a_\rho-B_h\rho$, with the exact $n=320$ coefficients
\begin{align*}
 A_h&=\frac{1241317}{9680000}-\frac{679}{251200}
                    -\frac{77392}{3768000}-\frac{97}{38400},\\
 B_h&=\frac{253869}{2420000}-\frac{35}{25120}
                    -\frac{2000}{150720}-\frac1{768}.
\end{align*}
Direct positive-denominator comparisons give
$5123/50000<A_h<1281/12500$ and $0<B_h<4447/50000$.
For the nearest and small-angle classes, respectively, the same floor
formulas give
\begin{gather*}
 F_n>\frac{27}{200}\left(\frac{959}{1000}a_\rho-\frac\rho2\right)
       >\frac{11}{100}(a_\rho-\rho/2),\\
 F_s=\frac{35}{286}\left(\frac{157}{160}\right)^2(a_\rho-\rho/2)
       >\frac{11}{100}(a_\rho-\rho/2).
\end{gather*}
The first comparison uses
$2[n/(n+1)](1/4-1/n)=158/321$ and
$\pi+\rho\leq51/14$.  The second-step difference for $F_n$ is
$3893a_\rho/200000-\rho/80>0$.  The common lower bound for these
two classes exceeds $L_h$ by
\[
 \frac{377}{50000}+\frac{1697}{50000}\rho
                                    -\frac{94}{50000}\rho^2>0.
\]
The three classes exhaust each eligible projection, and the original
floor formulas improve with $n$ at the same fixed radius.  Every
noncorner odd center has at least one eligible interior index; only
that projection is required.

Combine \eqref{eq:argmax-central-pair} and
\eqref{eq:argmax-beta-quadratic} with
$u\geq47(x^2+y^2)/192$.  The fixed tangent
$\sqrt x\leq(5/4)x+1/5$ follows from
$(\sqrt x-2/5)^2\geq0$.  On $1/8\leq\rho\leq1/2$,
\eqref{eq:argmax-base320} gives $(47/192)U_N>1143/400$.
Consequently
\begin{equation}
 \begin{gathered}
 S_n-\Lambda_c<C_4+\ell_1x+\ell_2y-d_1x^2-d_2y^2,
 \\
 \ell_1=\frac{12531}{20000},\quad\ell_2=\frac{11056}{20000},
 \quad d_1=\frac{1911}{800},\quad d_2=\frac{24539}{8000}.
 \end{gathered}
                                                        \label{eq:argmax-half-quadratic}
\end{equation}
For positive $d_1,d_2$, the maximum of the quadratic part on the entire
line $x+y=\rho$ is $-e\rho^2+f\rho+g$, where
\begin{equation}
 e=\frac{d_1d_2}{d_1+d_2},\qquad
 f=\frac{d_2\ell_1+d_1\ell_2}{d_1+d_2},\qquad
 g=\frac{(\ell_1-\ell_2)^2}{4(d_1+d_2)}.
                                                        \label{eq:argmax-radial-maximum}
\end{equation}
Enlarging the nonnegative segment to this line is a valid upper bound.
Here
\[
 e>13429/10000,\quad f<29714/50000,\quad g<1/4000,
\]
so
\[
 S_n-\Lambda_c<\frac{1379}{100000}
                  +\frac{29714}{50000}\rho-\frac{13429}{10000}\rho^2.
\]
Subtracting this from \eqref{eq:argmax-half-floor}, the resulting
quadratic exceeds $1/20000$ for every real $\rho$:
\[
 \frac{8867}{100000}-\frac{34161}{50000}\rho
                         +\frac{8233}{6250}\rho^2>\frac1{20000}.
\]
Completing the square reduces this assertion to
$34161^2=1166973921<1167373536=32\cdot8233\cdot4431$.

For $0\leq\rho\leq1/8$, drop the nonnegative radius term in $U_N$.
Both curvatures exceed two, whence
\[
 S_n-\Lambda_c<677/50000+(627/1000)\rho-\rho^2<77/1000.
\]
On the same interval $L_h(\rho)\geq L_h(1/8)>9/100$.
This proves \eqref{eq:argmax-half-near} on the whole closed half
diamond, including its axes and origin.

\subsection{The matched annulus and the wide primitive}

The new wide coefficients in \eqref{eq:argmax-coefficients} obey at
$n=320$
\[
 k_{Wx}<373019/640000<583/1000,\qquad
 k_{Wy}<59808113/102400000<117/200.
\]
Their inverse-square charges $152/n^2$ and $d_\beta$ are included.
Also $T-\bar e+j>2331/200$ and
$(11/48)(T-\bar e+j)>8/3$.  For $\rho\leq1$,
$u\geq11(x^2+y^2)/48$ and
$E_c\leq5x^{3/2}/12$ by Lemma~\ref{lem:argmax-central-pair}.
Using the fixed tangent $\sqrt x\leq x+1/4$ gives on the entire
nonnegative radius-one segment
\[
 S_n-\Lambda_c<C_4+\frac{11}{16}x+\frac{47}{80}y
                         -\frac94x^2-\frac83y^2.
\]
Formula~\eqref{eq:argmax-radial-maximum} now has
\[
 e=72/59>61/50,\quad f=3029/4720<257/400,\quad g=3/5900,
\]
and hence
\begin{equation}
 S_n-\Lambda_c<P(\rho):=
 \frac{141}{10000}+\frac{257}{400}\rho-\frac{61}{50}\rho^2.
                                                        \label{eq:argmax-wide-polynomial}
\end{equation}
$P$ decreases for $\rho\geq1/2$, and
$P(1/2)=607/20000$, $P(3/5)=-99/2500$.  This proves
\eqref{eq:argmax-wide}.

Only $1/2\leq\rho\leq3/5$ needs a defect.  The nearest and small-angle
classes give
\[
 F_n>\frac3{25}\left(\frac{19}{20}a_\rho-\frac\rho2\right)
       >\frac9{100}(a_\rho-\rho/2),\qquad
 F_s>\frac{11}{100}(a_\rho-\rho/2),
\]
whose common lower bound is at least $549/10000>17/500$.
For the middle class apply its actual extended guard
$\rho\leq21/40$, obtaining directly
\[
 F_h>\frac{5123}{50000}\left(1-\frac{(21/40)^2}{4}\right)
             -\frac{4447}{50000}\frac{21}{40}>\frac{17}{500}.
\]
For $21/40\leq\rho\leq3/5<\pi/4$, use the separately valid wide
floor $F_w=A_w a_\rho-B_w\rho$ with
\begin{align*}
 A_w&=\frac{2893891}{29040000}-\frac{679}{251200}
                    -\frac{75242}{3768000}-\frac{97}{38400},\\
 B_w&=\frac{3977281}{58080000}-\frac{35}{25120}
                    -\frac{15}{1256}-\frac1{768}.
\end{align*}
Here $37/500<A_w<3/40$ and $0<B_w<27/500$, so
\[
 F_w>\frac{37}{500}-\frac{27}{500}\rho-\frac3{160}\rho^2
       \geq\frac{697}{20000}>\frac{17}{500}.
\]
Every eligible projection is covered by these classes.  Subtracting
the defect from \eqref{eq:argmax-wide-polynomial} leaves a gap at least
$17/500-607/20000=73/20000>3/1000$, proving
\eqref{eq:argmax-annulus}.  This is a fresh use of the guarded floor
formulas; the half-diamond conclusion was not extended by assertion.
For all real $n\geq320$, $T-\bar e+j$ increases while the new
$k_{Wx},k_{Wy},C_4$ decrease, and the same-radius floor formulas improve.
Thus both annular conclusions hold throughout the stated integer tail.

\subsection{The full outer diamond}

For $1\leq\rho\leq\pi$ put
\[
 c_e=\cos(\rho/2),\quad c_d=\cos((x-y)/2),\quad a=c_ec_d,
 \quad u_1=47/384,\quad v=u-u_1,
\]
and fix
\[
 a_0=337/384,\quad w_0=721/384,\quad
 p_0=12/25,\quad P_0=24/25,\quad\eta=3/2.
\]
The following actual-phase bounds preserve the correlation between
phase radius and model loss:
\begin{equation}
 \begin{gathered}
 v\geq0,\quad a\leq a_0,\quad c_e+c_d\leq w_0,\qquad
 \max(p,q)\leq p_0+\eta\sqrt v,\\
 (q-p)_+\leq|p-q|\leq3\sqrt v,\qquad
 P=p+q\leq P_0+4v.
 \end{gathered}                                      \label{eq:argmax-outer-geometry}
\end{equation}
To verify them, write $c=\cos(1/2)<337/384$ and
$u_{\min}=1-c$.  Then $c_e\leq c$, $c_d\leq1$ and
$c_e+c_d\leq1+a$.  If $x\leq1/2$, then $p\leq\sin(1/2)<12/25$.
If $1/2\leq x\leq1$, write $h=x-1/2$ and use $y\geq1-x$ to obtain
\[
 u-u_{\min}\geq c(1-\cos h)\geq47ch^2/96,\qquad
 \sin x\leq\sin(1/2)+ch.
\]
Since $96c/47<9/4$, this proves the sine bound.  If $x\geq1$, then
$u\geq11/48$.  When $u\leq1/4$, $p\leq\sqrt3/2<9/10$ while
$p_0+\eta\sqrt v>93/100$; when $u\geq1/4$ the latter exceeds one.
The same trigonometric argument applies to $q$ without interchanging
the unequal finite meshes.  Further,
$|p-q|=2\sqrt{c_e^2-a^2}\leq2\sqrt{2v}$.
Finally $P\leq2\sqrt{2u-u^2}$ by Cauchy.  This last function is
concave; its value at $u_1$ is $\sqrt{33887}/192<24/25$ and its
derivative is $674/\sqrt{33887}<4$, proving its stated tangent bound.

We make the absolute-error envelope explicit.  Set
\[
 \epsilon=3\pi/(2m),\quad\ell=-\log\sin\epsilon,\quad
 D_a=D-10/n^2,\quad C=1+8a_0m/(3\pi),
\]
and define
\begin{align*}
 \Gamma_b&=p_0[1+\log(1+C/p_0)],\\
 \gamma_b&=1+\log(1+C/p_0)-C/(p_0+C),\\
 M_{A,b}&=p_0/\pi+4a_0/\pi^2,\\
 C_{B,b}&=6+(4+4a_0/\mathrm e)/\pi,\\
 V_{B,b}&=23+[3\Gamma_b+(2p_0+6P_0)\ell+26a_0+20]/\pi,\\
 G_b^{\rm env}&=7LM_{A,b}/m+V_{B,b}/(2m)\\
 &\hspace{8mm}+(4/m)[7+(P_0/\pi)(5+2\log(2m))]+2C_{B,b}/m,\\
 D^{\rm env}&=a_0D_a+10/n^2,\\
 R_6^{\rm env}&=\frac{21+21a_0+2L+[(1+48p_0)H+24p_0]/\pi}{n}
                                    +\frac{317+5/12}{n^2},\\
 D_m&=4w_0/\pi+a_0b_0\log\pi,\qquad Z_m=D_m+a_0b_0/n,\\
 \mathcal M^{\rm env}&=\frac{a_0b_0(2L+2h_r)}n
  +\frac{D_m+a_0b_0(L+1/n)}{4n^2-1}+\frac{a_0b_0}{(2n-1)^2}\\
 &\hspace{8mm}+\frac{Z_m+a_0b_0(\log(n+1/2)+2)}{4(n-1/2)^2},\\
 C_b&=D^{\rm env}+G_b^{\rm env}+R_6^{\rm env}+\mathcal M^{\rm env},\\
 C_s&=\eta\left\{\frac{7L}{\pi m}
      +\frac{3\gamma_b+2\ell}{2\pi m}+\frac{48H+24}{\pi n}\right\}
                    +3\left\{\frac{\log2+4\log\pi}{\pi m}+W\right\},\\
 C_l&=\frac{12\ell+16(5+2\log(2m))}{\pi m}.
\end{align*}
Applied first at an actual phase, the full absolute row theorem and
\eqref{eq:argmax-outer-geometry} give
\begin{equation}
 |S_n-M_\infty|\leq C_b+P_0K_D+C_s\sqrt v+(C_l+4K_D)v.
                                                        \label{eq:argmax-outer-error}
\end{equation}
For example the only nonlinear positive allowance is
$\Gamma(p,a)=p[1+\log(1+(1+8am/(3\pi))/p)]$, with zero limit at
$p=0$.  It increases in $a$ and is concave in $p$; its tangent at
$(p_0,a_0)$ is $\Gamma_b+\gamma_b(p-p_0)$.
This explains the $\gamma_b\sqrt v$ term.  All other terms follow by
positive majorization of the full ordinary, six retained, outer,
mesh, and scalar-log allowances.  Thus no fictitious tuple of envelope
parameters is substituted into a physical theorem.

Let $j_*=23/250$, $t_*=221/500$.  The global model master and the
whole one-sided origin cost yield
\[
 \Lambda_c-S_n\geq G+Kv-A\sqrt v,
\]
where
\begin{align}
 G&=u_1(T+j_*)-t_*p_0-C_b-P_0K_D-\bar e,\nonumber\\
 K&=T+j_*-C_l-4K_D,\qquad A=\eta t_*+C_s.
                                                        \label{eq:argmax-outer-gate}
\end{align}
At $n=320$, the rational logarithm bounds above give
\begin{gather*}
 K_D<1/6,\quad W<3/125,\quad\ell<13/3,\\
 \Gamma_b<18/5,\quad\gamma_b<15/2,
 \quad M_{A,b}<13/25,\quad C_{B,b}<8.
\end{gather*}
For $\ell$, use $\epsilon<1/7$ and $\sin\epsilon>23/(5m)$.
The inequality $\log(5m/23)<13/3$ follows from
\[
 \exp(13/3)>(8/3)^4(1+1/3+1/18)=51200/729>1605/23.
\]
The complete positive envelopes then obey
\[
 D^{\rm env}<3/25,\quad G_b^{\rm env}<9/25,
 \quad R_6^{\rm env}<17/50,\quad\mathcal M^{\rm env}<2/125,
 \quad C_b<209/250.
\]
An inspectable bound for the retained-row numerator is
\begin{align*}
 R_6^{\rm env}&<\frac{21+21(22/25)+2(577/100)}{320}\\
 &\quad+\frac{(601/25)(677/100)+288/25}{(31/10)320}
       +\frac{318}{320^2}<\frac{17}{50}.
\end{align*}
For $C_s$, the three terms in its $\eta$ bracket are bounded,
in order, by $41/1000,2/125,353/1000$; its remaining bracket is
less than $3/500+3/125$.  Therefore
$C_s<141/200<71/100$.  Also $C_l<7/20$.
With $\bar e<5/32$, these estimates give
\[
 G>\frac{195191}{2400000}>\frac3{40},\qquad K>10,
 \qquad A<\frac75.
\]
In particular $K$ is positive before division.  Completing the square
in $\sqrt v$ now yields
\[
 \Lambda_c-S_n\geq G-\frac{A^2}{4K}
       >\frac3{40}-\frac{(7/5)^2}{40}=\frac{13}{500},
\]
which is \eqref{eq:argmax-outer}.

The transfer to all real $n\geq320$ can be checked term by term.
$T$ increases, while $K_D,\bar e,C_b,C_s,C_l$ decrease.  In the
entropy terms, writing $C=1+cm$ gives
$m\Gamma_b'<p_0\leq\Gamma_b$ and
$m\gamma_b'<1<\gamma_b$, so their quotients by $m$ decrease.
Also $m\ell'=\epsilon\cot\epsilon\leq1<\ell$, hence $\ell/m$
decreases.  The mesh quotients have positive linear-log numerators and
quadratic denominators; their derivatives are negative on $n\geq32$.
It follows that $G,K$ increase and $A$ decreases, preserving the whole
square-completion gate.  No monotonicity of the actual $S_n$ is used.

\subsection{A strict patch at each actual corner}

We first state the interpolation fact that also supports the numerical
corner branch.  Suppose $M$ bounds every complete signed cardinal sum
globally, and put $k=n/m$.  If the three actual losses at
$(A,B)=(\sqrt{2U},0),(0,\sqrt{2V}),(\sqrt{2U},\sqrt{2V})$ are at least
$cU,cV,c(U+V)$, respectively, then
\begin{equation}
 \lambda_n(-\cos(A/n),-\cos(B/m))
 \leq\Lambda_c-\left[c-\frac M6\max(U,k^4V)\right]\frac{A^2+B^2}{2}
                                                        \label{eq:argmax-squared-peano}
\end{equation}
on the whole rectangle, provided the bracket is positive.
Indeed apply the common nonnegative interpolation kernels to the
smooth signed cardinal vector in $s=A^2/2,t=B^2/2$.  Its second
derivative bound is $M/3$ in $s$, $k^4M/3$ in $t$.  The Peano mass
therefore gives
\[
 H(s,t)\leq I H(s,t)+\frac M6\{s(U-s)+k^4t(V-t)\}.
\]
The vertex losses and positive bilinear weights give
$I H\leq\Lambda_c-c(s+t)$ exactly; the remainder is at most
$(M/6)\max(U,k^4V)(s+t)$.  This proves
\eqref{eq:argmax-squared-peano}, including both closed axes and the
origin.  The signed-cardinal second derivative bound follows by the
fourth-derivative Bernstein inequality and the integral identity
\[
 \frac{d^2}{ds^2}f(\sqrt{2s})
 =\int_0^1\int_0^1 t^2 f^{(4)}(rt\sqrt{2s})\,dr\,dt
\]
for an even smooth function $f$; its kernel has mass $1/3$.

For the analytic patch put $h=\pi/4$, $q_0=7/48$, and take
$U=V=h^2/2=\pi^2/32$.  The exact axis deficit
$(1-\cos h)/2>q_0$.  The old complete wide finite allowance, legal at
$\rho=h$, uses
\[
 k_{Wy}^{\rm old}
 =\frac{(19/2)H+15L+265/3}{n}+\frac{179+17L}{n^2}.
\]
It gives a lower bound for both actual axis losses,
\[
 \delta_A=q_0(T-\bar e+j)-\frac14-C_4-hk_{Wy}^{\rm old}.
\]
Here $E_c(h,0)<1/4$ follows from $\kappa_4<9/25$,
$a>17/20$ and $p^{3/2}<3/5$; $E_c(0,h)=-B_\beta\leq0$.
At the diagonal $(h,h)$ the completed-model theorem has its separately
proved closed $\rho\leq\pi/2$ domain.  Use only the full absolute
finite error there, not a relative radius-one estimate.  Since
$a=p=q=1/\sqrt2$, $P<3/2$, and $E_c<4/15$, the actual diagonal
loss exceeds
\[
 \delta_D=2q_0(T+j)-\frac4{15}-\overline C_E-\frac32K_D-\bar e.
\]

At the single base $256=2^8$, the positive atanh series gives
$277/50<L<111/20$ and $H<131/20$.  Direct substitution of all
positive terms gives
\[
 \begin{gathered}
 T>107/10,\quad D<1/6,\quad G_g<5/8,\quad
 (64+16L)/256+318/256^2<61/100,\\
 \mathcal M<11/500,\quad W<7/250,\quad
 \overline C_E<3/2,\quad\bar e<1/5,\quad K_D<1/5,\\
 C_4<9/500,\quad k_{Wy}^{\rm old}<93/100,\quad M(256)<26.
 \end{gathered}
\]
The norm bound is used only at the base; $M(n)$ grows with $n$.
For an explicit check on the largest collected error, at this base
$g<7$, $\log m<139/25$, $\log(2m)<63/10$, so
\[
 G_g<\frac{2331/80+111/2+2190/31}{257}<5/8.
\]
The complete mesh has its same three tails below $14/n^2$, and the
leading term is below $5371/256000$, proving its displayed bound.
The norm follows from
$M(256)<4218/197+9/4+3/2+2/5<26$.

Choose $r_U=2/125>\pi^4/6144$ and consider the combined gates
$F_A=\delta_A-r_UM(n)$, $F_D=\delta_D-2r_UM(n)$.
The base bounds imply
\[
 \delta_A>2293/4200,\quad F_A>2729/21000>1/8,
 \qquad \delta_D>881/1000,\quad F_D>49/1000.
\]
Both gates increase for all real $n\geq256$.  To see this without
treating the base norm as uniform, expand their leading polynomial:
for $F_A$ it is
\[
 \frac{4(q_0-r_U)L^2-[3q_0+(79/5)r_U]L}{\pi^2},
\]
and for $F_D$ twice that expression.  Its derivative in $L$ is
positive for $L\geq1$, because $8(q_0-r_U)=779/750$ and
$3q_0+(79/5)r_U=6903/10000$.  All remaining nonconstant costs are
positive decreasing combinations of
$\bar e,C_4,k_{Wy}^{\rm old},\overline C_E,K_D$.

With $\delta=\min(\delta_A,\delta_D/2)$ it follows that
\[
 \delta-M(n)\pi^4/6144>49/2000.
\]
In \eqref{eq:argmax-squared-peano}, take $c=32\delta/\pi^2$ and
use $k\leq1$ only in the nonnegative remainder.  The coefficient of
$A^2+B^2$ is greater than
$(16/\pi^2)(49/2000)>49/1250>1/30$.
This proves \eqref{eq:argmax-corner-loss} for every $n\geq256$.
The node-preserving parity reflection carries the same proof to the
other actual corner, without exchanging $n,m$.

\subsection{Exact phase-certificate predicates}
\label{app:argmax-certificate-predicates}

We spell out the finite implication used in
Proposition~\ref{prop:certificate-phase-soundness}.  Let $p_-\leq\pi\leq p_+$
be positive rationals, let $L_c\leq\Lambda_c(n)$, and put $k=n/m$.
The cover is a complete binary midpoint tree rooted at $[0,1]^2$.
Both children are closed, and every split has positive rational width.
For a leaf $B=[a_l,a_h]\times[b_l,b_h]$, discard it only when
$a_l+b_l>1$.  Equality is a retained tangent intersection with
$\mathcal T$.  Complete tree syntax, both subtrees, and complete token
consumption prove the closed cover by induction.

For vertex upper endpoints $W_z\geq S_n^t(z)$ and a genuine global
norm $M$, put
\begin{align}
 \theta_B&=\frac{p_+^2}{8}
                   [(a_h-a_l)^2+k^2(b_h-b_l)^2],\nonumber\\
 q_B(a,b)&=I_BW(a,b)+\frac{Mp_+^2}{2}
  \{(a-a_l)(a_h-a)+k^2(b-b_l)(b_h-b)\},\nonumber\\
 V_B&=\max_{z\ {\rm vertex\ of}\ B}W_z,\qquad U_B=V_B+\theta_BM.
                                                        \label{eq:argmax-certificate-peano}
\end{align}
Then $S_n^t\leq q_B\leq U_B$ on the entire rectangle.  The
$\pi^2$ factor is essential: the frequencies are at most $1,k$ in
$\alpha,\beta$, hence at most $\pi,\pi k$ in $a,b$.
This is the common signed-vector Peano argument, not differentiation
of the nonsmooth absolute sum.

Let $P_B=B\cap\mathcal T$, $\ell_B=a_l+b_l$, and
$h_B=\min(a_h+b_h,1)$.  These are the exact extrema of $a+b$ on
$P_B$.  The actual near and far parts are contained, respectively, in
the rational proxies
\[
 Q_N=P_B\cap\{a+b\leq1/(2p_-)\},\qquad
 Q_F=P_B\cap\{a+b\geq1/(2p_+)\}.
\]
The proxies may overlap and retain tangencies.  A noncorner floor
$\delta_B>0$ is proved on the signed radius
\[
 r_B=\min\{1/2,p_+(a_h+b_h)\}.
\]
It is applied only at an actual near point: there
$\rho=\pi(a+b)\leq1/2$ and $\rho\leq p_+(a_h+b_h)$.
Membership in a near proxy alone is not a defect-domain premise.

Each numerical leaf must pass one of the following exact rules.
\begin{enumerate}
 \item A far leaf has a validated upper $V\geq\max_{P_B}q_B$ with
 $V<L_c$.
 \item A near leaf has $2p_+h_B\leq1$ and an upper
 $V\geq\max_{P_B}q_B$ with $V-\delta_B<L_c$.
 \item A split leaf has validated uppers $V_N,V_F$ on the two
 nonempty proxies and satisfies $V_N-\delta_B<L_c$ and $V_F<L_c$
 for each nonempty corresponding proxy.
\end{enumerate}
An empty proxy makes only its actual branch vacuous.  Taking all
uppers equal to $U_B$ is a valid conservative choice.  The near rule
and near half of the split rule concern noncorner centers; at an
actual corner the direct punctured-corner witness is used instead.
Every comparison is strict.  An analytic far leaf is allowed only
under its entire closed-box guard, as proved below.

The global norm has its own finite witness.  For a wholly numerical
tree, let $M_* =\max_{\mathcal T}S_n^t$, which is finite by continuity.
The even-lattice period and separate phase sign symmetries imply that
this is the global real-phase norm.  Applying the Peano bound with
the true $M_*$ to a retained leaf containing a maximizer gives, if
every $\theta_B<1$,
\begin{equation}
 M_*\leq\max_B\frac{V_B}{1-\theta_B}.
                                                        \label{eq:argmax-tree-norm}
\end{equation}
This absorption uses every final nonempty numerical leaf, including
tangencies.  It does not assert equality between a sampled and a
continuous maximum.

Alternatively retain the quarter-grid points
$E=\{(i/4,j/4):i,j\geq0,\ i+j\leq4\}$ and their upper maximum $V_E$.
The ten closed quarter squares with lower-left indices $i+j\leq3$
cover $\mathcal T$.  Their outside vertices fold into $E$ by
$(a,b)\mapsto(1-a,1-b)$, an even-lattice/sign symmetry.  The Peano
remainder is at most $\pi^2(1+k^2)M_*/128\leq\pi^2M_*/64$.
Therefore
\begin{equation}
 M_*\leq\frac{V_E}{1-\pi^2/64}<\frac{784}{663}V_E.
                                                        \label{eq:argmax-grid-norm}
\end{equation}
The $15$ triangular grid points have $13$ canonical keys, since only
on the hypotenuse may the symmetry $(a,b)\leftrightarrow(b,a)$ be used.
Interior coordinate exchange is not a symmetry used by this argument.
All $13$ grid bounds are retained when this norm is used.

A tree containing an analytic far leaf must use the grid norm.
Such a leaf has no numerical upper for the global norm absorption in
\eqref{eq:argmax-tree-norm}; the remaining numerical leaves need not
contain a maximizer.  For a wholly numerical tree the final norm may
be the smaller valid rounded grid or tree bound.  On any change of
norm, all final Peano, leaf, and norm-dependent corner predicates are
recomputed.  No discarded sample or discovery-stage norm remains a
hidden premise.

The numerical actual-corner witness uses the three points
$(\pi/4,0),(0,3\pi/16),(\pi/4,3\pi/16)$ in inward phase coordinates.
With their actual upper endpoints $W_x,W_y,W_{xy}$, put
\[
 u'=1/32,\quad v'=9/512,\quad
 q=\min\left\{\frac{L_c-W_x}{u'},\frac{L_c-W_y}{v'},
                         \frac{L_c-W_{xy}}{u'+v'}\right\}.
\]
The exact strict predicate is
\begin{equation}
 q>\frac{Mp_+^4}{6}\max\{u',k^4v'\}.             \label{eq:argmax-finite-corner-test}
\end{equation}
In \eqref{eq:argmax-squared-peano}, take
$U=\pi^2u'$, $V=\pi^2v'$, and $c=q/\pi^2$.
The resulting radial loss is at least
\[
 \frac{q-(Mp_+^4/6)\max(u',k^4v')}{2p_+^2}(A^2+B^2)>0
\]
off the origin.  The true $\Lambda_c$ is the origin interpolation
entry; the numerical lower $L_c$ is used only in the three losses.
Both phase endpoints exceed $1/2$, so the rectangle covers the whole
inward half diamond.  Its physical angle-turn inputs are
$(1/(4n),0),(0,3/(16m)),(1/(4n),3/(16m))$.
They are not unscaled phase turns.

\subsection{The analytic far-leaf guard}

For completeness, the hybrid certificates use the following separate
regional consequence of the absolute row bound:
\begin{equation}
 S_n<\Lambda_c(n)-79/2000<\Lambda_c(n)-3/100,
 \qquad n\geq192,\quad \pi/2\leq\rho\leq\pi.
                                                        \label{eq:argmax-far192}
\end{equation}
Here is its quantitative reduction.  For $0<u\leq1$ set
\[
 A_u=1-u,\quad C_u=2-u,\quad
 F_u=\begin{cases}2\sqrt{u(1-u)},&u\leq1/2,\\1,&u\geq1/2,\end{cases}
 \qquad Q_u=\min(2,\sqrt{8u}).
\]
Actual phases satisfy $a=A_u$, $c_e+c_d\leq C_u$,
$p,q,(q-p)_+\leq F_u$, and $P\leq Q_u$.
Let $\widehat C_E(n,u)$ denote the explicit positive phase allowance
of Lemma~\ref{lem:phase-row-global} after replacing these quantities
by their displayed upper bounds.  This is a majorization of the
allowance at an actual phase, not an application of its physical
theorem to independent artificial parameters.  The quotient
$\widehat C_E(n,u)/u$ decreases: all nonentropy terms are positive
combinations of $1/u,A_u/u,C_u/u,F_u/u,Q_u/u$ with decreasing factors.
For entropy, the increasing homogeneous function
$\Psi(r,s)=r[1+\log(1+s/r)]$ gives
\[
 \widehat\Gamma/u
 =\Psi\left(F_u/u,[1+8A_um/(3\pi)]/u\right),
\]
whose two arguments decrease.  The zero limits and the two branches
at $u=1/2$ are continuous.

On $\rho\geq\pi/2$, $a\leq\cos(\pi/4)<17/24$, so $u>u_0=7/24$.
At that fixed anchor,
$A_u=17/24$, $C_u=41/24$, $F_u=\sqrt{119}/12<11/12$,
and $Q_u=\sqrt{7/3}<31/20$.  At $n=192=3\cdot2^6$, the atanh
series brackets $\log3$ between $7379/6720$ and that sum plus
$1/1728$.  Together with the preceding $\log2$ bracket this gives
$21/4<L<421/80$, $H<63/10$, $\log m<53/10$, $\log(2m)<6$.
Substitution into every term of the positive allowance gives
\[
 \begin{gathered}
 T>19/2,\quad D<43/200,\quad\mathcal M<7/250,
 \quad\bar e<1/4,\quad K_D<13/50,\\
 \widehat D<4/25,\quad\widehat G<7/10,
 \quad\widehat C_6<31/40,\quad\widehat{\mathcal M}<1/50,
 \quad F_uW<1/30,
 \end{gathered}
\]
and hence $\widehat C_E(192,u_0)<1013/600<17/10$.
These bounds retain all ordinary, six-row, mesh, outer, and scalar-log
terms.  For example $\ell<4$, $\widehat\Gamma<13/2$,
$\widehat M_A<3/5$, $\widehat V_B<57$, and $\widehat C_B<8$;
the five numerator bounds for $\widehat G$ are
$23,29,62,2,16$, summing to less than $132$ before division by $193$.
The retained-row numerator is below $146$, with remainder
$318/192^2$; all three weighted mesh tails total less than $14/192^2$.

The global model master now gives, at the anchor,
\begin{align*}
 \Lambda_c-S_n&>
 u_0(19/2+23/250)-(221/500)(11/12)
                   -17/10-(13/50)(31/20)-1/4\\
 &=79/2000.
\end{align*}
The quotient of this lower bound as a function of $u$ increases,
by the quotient property just proved and the decrease of $F_u/u,Q_u/u$.
Thus it stays positive for every $u\geq u_0$.  At each fixed actual
phase the original absolute allowance, $K_D$, and $\bar e$ decrease
with real $n$, while $T$ increases.  This transfers the comparison to
every integer $n\geq192$ and proves \eqref{eq:argmax-far192}.

An explicit analytic leaf must have $n\geq192$ and its raw rational
box must satisfy $0\leq a_l<a_h\leq1$, $0\leq b_l<b_h\leq1$,
and $1/2\leq a_l+b_l\leq1$.  Its entire intersection with
$\mathcal T$ then has $\pi/2\leq\rho\leq\pi$, including closed seams.
It is compared with the true symbolic $\Lambda_c$, and supplies no
numerical norm-support point.  These guards explain both the strict
analytic far exclusion and the compulsory grid norm in a hybrid tree.

\subsection{Extent of the finite verification}
\label{app:argmax-verification-record}

The inherited finite range is $2,\ldots,40$.  Degrees $33$--$40$ use
the corrected-cardinal Bernstein constructor, a separately authored
C++/GMP checker, complete prefix partitions, and two strict punctured
corner leaves.  The essential certificates contain, respectively,
$284,288,298,306,316,320,330,350$ bound leaves and two corner leaves.
Earlier exact polynomial cases cover $2$--$7$, and complete exact or
outward Bernstein covers cover $8$--$32$.

The subsequent returned phase worklist is exactly $41,\ldots,319$
together with the control degree $511$.  Its $280$ primary/fresh pairs
include all $273$ residual degrees
\[
 \{43,\ldots,62\}\cup\{65,\ldots,158\}\cup\{161,\ldots,319\}
\]
and the seven controls $41,42,63,64,159,160,511$.  Both retained
certificates in every pair were reconstructed locally, for a total of
$560$ exact rational proof reconstructions.  Every degree identity,
whole cover, final norm, floor, corner or far guard, strict inequality,
source link, final account, and recorded zero exit was checked.
This is a setwise degree verification, not merely a count of receipts.

The final hybrid scientific freeze has SHA-256
\begin{center}\ttfamily\footnotesize
 0d9ce0a5dbb7fb37cf1a05f07040b72d\\
 046b61b8d36b6a5a8bf137eb7d9914eb.
\end{center}
The production source commit is
\texttt{996fedcad4ba79a28e45a3fe3678adc12f1e5358}; the runtime manifest
and all transitive source pins are in the machine-readable evidence.
All rational predicates use lossless conversion of canonical finite
binary64 hexadecimal endpoints.  The native inclusion uses the frozen
corrected finite row evaluator and its Arb/FLINT/MPFR/GMP outward
arithmetic, including polynomial collision values.  Exact cover checking
does not independently prove the native arithmetic or compiler correct.

On the production machine, primary and retained-witness fresh passes
ran separately for all pairs, and the released collector validated
their full histories.  The compact return retains all essential proofs
and their source/account/exit links.  Local independent checking of
the unsampled cases reconstructs the rational implications and checks
these links, with the outward endpoints as native witnesses; it does
not inspect full unsent native histories.

The persisted draw seed is
\begin{center}\ttfamily\footnotesize
 7c79b8e401d9a32f8233f8371c441cd2\\
 8d82dedd60e46e7c73d211ab099d01d2.
\end{center}
Sampling two degrees from the sorted $273$-degree residual population
with the recorded Python random routine gives the ordered pair
$199,238$.  Both original full closures were inspected locally.
Exactly two new native verify-mode executions then used their retained
primary witnesses and fresh exclusive output paths, with no replacement
sample or retry.  Both exited zero with complete final accounts; each
used $45$ common-primitive calls and three actual-corner calls.  A
distinct reviewer validated the fresh results and their source/runtime
closures.  Both samples are below $256$; previous degree-$511$ checks
and the mathematical and source reviews support the analytic-corner
branch.  Fresh execution reuses the same arithmetic algorithm and is
not an independent arithmetic implementation.

The actual compact archive has $36{,}045{,}486$ bytes and SHA-256
\begin{center}\ttfamily\footnotesize
 b00fd57f85ce6d96a6c1fa01e2307abd\\
 815088bb8fcb28b1587657ef142b109a.
\end{center}
All essential unsampled proofs remain in that compact return; the two
sampled complete histories and small verification records are also
retained locally.  Full unsampled histories remain on the production
machine under the explicitly accepted selective-retention policy.
External durable retrieval of the new compact return has not been
verified.  Source pins, compatible-runtime identities, degree lists,
and commands permit regeneration of valid witnesses; they do not
guarantee recovery of every historical timing, path, or original byte.
These limits are part of the evidence description, not an inference
that the two random checks prove unsampled numerical premises.
